\documentclass[11pt, a4paper]{article}
\usepackage[a4paper, margin=2.5cm]{geometry}
\usepackage{fontspec}

% Configuración limpia de Babel para LuaLaTeX / XeLaTeX
\usepackage[bidi=basic, provide=*]{babel}
\babelprovide[import, main, onchar=ids fonts]{spanish}
\babelprovide[import, onchar=ids fonts]{english}

\babelfont{rm}{Noto Sans}

% Microtipografía y tolerancia ante cajas ajustadas
\usepackage{microtype}
\emergencystretch=2.5em

\usepackage{enumitem}
\setlist[itemize]{label=-}
\usepackage{amsmath, amssymb}
\usepackage{booktabs}
\usepackage{tabularx}
\usepackage{hyperref}

\title{\textbf{Programa de Posgrado: Geometría Lineal}\\ \large Curso de Estudios Dirigidos --- Maestría en Computación y Matemática}
\author{David Duarte G.}
\date{Modalidad 100\% Virtual (Moodle)}

\begin{document}
	
	\maketitle
	
	\section*{Ficha Técnica del Curso}
	
	\begin{center}
		\begin{tabularx}{\textwidth}{lX}
			\toprule
			\textbf{Asignatura:} & Geometría Lineal \\
			\textbf{Nivel:} & Maestría en Computación y Matemática \\
			\textbf{Modalidad:} & 100\% en Línea (Plataforma Moodle) \\
			\textbf{Régimen:} & Estudios Dirigidos \\
			\textbf{Carácter:} & Teórico-Práctico / Computacional \\
			
			\bottomrule
		\end{tabularx}
	\end{center}
	
	\section{Fundamentación y Descripción General}
	
	La geometría moderna constituye el lenguaje formal indispensable para articular el análisis matemático riguroso con el diseño de algoritmos espaciales en computación visual, robótica móvil, procesamiento digital de geometrías y simulación física. Tradicionalmente, la geometría diferencial y el álgebra lineal se enseñan desde perspectivas analíticas abstractas desconectadas de sus retos numéricos.
	
	Este curso de posgrado en la modalidad de estudios dirigidos cubre desde la fundamentación axiomática de los espacios afines y euclidianos hasta la geometría diferencial de curvas y superficies en $\mathbb{R}^3$, integrando en cada unidad derivaciones simbólicas exactas con implementaciones vectorizadas reproducibles en Python y Julia. A través del entorno virtual Moodle, el estudiante interactúa con cuadernos interactivos, foros de discusión crítica sobre estabilidad numérica y talleres de revisión por pares.
	
	\section{Objetivos de Aprendizaje}
	
	\begin{itemize}
		\item Formular algebraicamente las estructuras afines y euclidianas, diferenciando formalmente puntos de vectores libres y aplicando coordenadas baricéntricas en algoritmos de computación gráfica.
		\item Describir cuantitativamente la geometría intrínseca y extrínseca de curvas espaciales mediante el triedro de Frenet-Serret, modelando trayectorias cinemáticas continuas para robots y maquinaria CNC.
		\item Calcular y analizar la primera y segunda formas fundamentales de superficies regulares, vinculando las curvaturas principales, gaussiana y media con la segmentación de imágenes 3D y la deformación elástica de mallas.
		\item Aplicar el análisis tensorial y el Teorema Egregium de Gauss a la solución de ecuaciones geodésicas en variedades bidimensionales y problemas de optimización sobre superficies curvas.
	\end{itemize}
	
	\section{Unidades Temáticas Detalladas}
	
	\subsection{Unidad 1: Espacios Afines, Euclidianos y Álgebra Matricial Geométrica}
	
	\subsubsection*{Contenidos Analíticos}
	\begin{itemize}
		\item Estructura axiomática de espacio afín: acción transitiva y libre de un espacio vectorial sobre un conjunto de puntos. Regla de Chasles.
		\item Coordenadas baricéntricas, combinaciones convexas y símplices afines.
		\item Subespacios y variedades afines: rectas, planos e hiperplanos.
		\item El grupo afín $\text{Aff}(n)$ e isometrías euclidianas.
		\item Representación en coordenadas homogéneas y transformaciones matriciales por bloques.
		\item Descomposición en Valores Singulares (SVD) orientada a transformaciones geométricas, rotaciones óptimas y el problema de Procrustes ortogonal.
	\end{itemize}
	
	\subsubsection*{Aplicaciones Computacionales en el Mundo Real}
	\begin{itemize}
		\item \textbf{Robótica y Cinemática:} Encadenamiento de marcos afines locales mediante matrices de Denavit-Hartenberg para manipuladores seriales.
		\item \textbf{Computación Gráfica:} Rasterización de triángulos mediante interpolación baricéntrica acelerada por GPU.
		\item \textbf{Visión Artificial:} Estimación de matrices de transformación rígida y registro de nubes de puntos 3D (algoritmo Iterative Closest Point, ICP) resolviendo el problema de Procrustes vía SVD.
	\end{itemize}
	
	\subsubsection*{Actividades y Evaluaciones en Moodle}
	\begin{itemize}
		\item \textbf{Foro Académico 1:} Discusión formal sobre la preservación de invariantes afines frente a isometrías euclidianas en visión por computadora.
		\item \textbf{Laboratorio Computacional 1 (Jupyter Notebook):} Implementación desde cero de un registrador de nubes de puntos 3D mediante descomposición SVD.
	\end{itemize}
	
	\subsection{Unidad 2: Teoría Diferencial de Curvas y Cinemática de Trayectorias}
	
	\subsubsection*{Contenidos Analíticos}
	\begin{itemize}
		\item Concepto de curva: representaciones paramétricas regulares, curvas de clase $C^m$, representaciones implícitas y proyecciones ortogonales.
		\item Longitud de arco y parametrización natural por longitud de arco.
		\item Geometría local: vector tangente unitario, recta tangente, curvatura ($\kappa$), vector unitario normal principal, plano osculador, vector binormal y torsión ($\tau$).
		\item El triedro móvil y las Fórmulas de Frenet-Serret.
		\item Teorema fundamental de existencia y unicidad de curvas en $\mathbb{R}^3$. Ecuaciones intrínsecas y representación canónica local.
		\item Involutas, evolutas, teoría del contacto, circunferencias y curvas osculatrices.
	\end{itemize}
	
	\subsubsection*{Aplicaciones Computacionales en el Mundo Real}
	\begin{itemize}
		\item \textbf{Planificación de Movimiento en Vehículos Autónomos:} Generación de trayectorias sin discontinuidades de aceleración lateral mediante clotoides (curvas de Euler) y B-splines de clase $C^2$.
		\item \textbf{Mecanizado CNC e Impresión 3D:} Control dinámico del cabezal de corte en función de la curvatura local para minimizar vibraciones y desgaste mecánico.
		\item \textbf{Diseño Mecánico Asistido por Computadora (CAD/CAM):} Perfiles de dientes de engranajes conjugados construidos mediante involutas de círculo.
		\item \textbf{Cinematografía Virtual:} Interpolación de orientaciones de cámaras a lo largo de curvas con torsión controlada para prevenir inversiones bruscas de perspectiva (\emph{gimbal lock}).
	\end{itemize}
	
	\subsubsection*{Actividades y Evaluaciones en Moodle}
	\begin{itemize}
		\item \textbf{Laboratorio Computacional 2 (Jupyter Notebook):} Generador e interpolador numérico del triedro de Frenet sobre trayectorias 3D con animación vectorial continua.
		\item \textbf{Taller Moodle (Evaluación por Pares):} Deducción analítica y simulación computacional de la evoluta de curvas cónicas.
	\end{itemize}
	
	\subsection{Unidad 3: Teoría Diferencial de Superficies y Análisis Tensorial}
	
	\subsubsection*{Contenidos Analíticos}
	\begin{itemize}
		\item Concepto de superficie: parametrizaciones regulares, cartas locales, superficies simples. Plano tangente y recta normal. Propiedades topológicas básicas.
		\item Primera Forma Fundamental ($I$): tensor métrico, cálculo de longitud de arco sobre superficies, áreas y medidas angulares.
		\item Segunda Forma Fundamental ($II$): curvatura normal, curvaturas principales ($\kappa_1, \kappa_2$) y direcciones principales.
		\item Curvatura Gaussiana ($K = \kappa_1 \kappa_2$) y Curvatura Media ($H = \frac{\kappa_1 + \kappa_2}{2}$). Líneas de curvatura, líneas asintóticas y familias conjugadas de curvas.
		\item Ecuaciones de Gauss-Weingarten, relaciones de compatibilidad de Codazzi-Mainardi y el \emph{Theorema Egregium} de Gauss.
		\item Teorema fundamental de la teoría de superficies.
		\item Introducción al análisis tensorial: tensores covariantes y contravariantes, álgebra tensorial y formulación tensorial de la curvatura.
	\end{itemize}
	
	\subsubsection*{Aplicaciones Computacionales en el Mundo Real}
	\begin{itemize}
		\item \textbf{Imágenes Médicas y Biometría 3D:} Clasificación de patologías superficiales y segmentación de órganos a partir del mapa de signos de $K$ y $H$ (puntos elípticos, hiperbólicos, parabólicos y planos).
		\item \textbf{Mapeo de Texturas y Despliegue UV:} Minimización de la distorsión métrica en mallas tridimensionales preservando la primera forma fundamental.
		\item \textbf{Procesamiento Digital de Geometrías:} Algoritmos de suavizado mediante flujo de curvatura media (\emph{Mean Curvature Flow}) para filtrado de ruido en escaneos láser.
		\item \textbf{Cálculo de Estructuras por Elementos Finitos (FEM):} Análisis de membranas delgadas y cáscaras arquitectónicas mediante tensores de esfuerzo y deformación métrica.
	\end{itemize}
	
	\subsubsection*{Actividades y Evaluaciones en Moodle}
	\begin{itemize}
		\item \textbf{Laboratorio Computacional 3 (Jupyter Notebook):} Cálculo de las formas fundamentales y visualización del mapa térmico de curvatura gaussiana sobre archivos de mallas poligonales (.OBJ).
		\item \textbf{Cuestionario Numérico Moodle (STACK):} Evaluación de cálculo simbólico y numérico de coeficientes métricos $E, F, G$ y $e, f, g$.
	\end{itemize}
	
	\subsection{Unidad 4: Geometría Intrínseca, Geodésicas y Tópicos Avanzados}
	
	\subsubsection*{Contenidos Analíticos}
	\begin{itemize}
		\item Aplicaciones diferenciables entre superficies: isometrías locales y globales, aplicaciones conformes y equiareales.
		\item Geometría intrínseca: conexión de Levi-Civita, derivada covariante y curvatura geodésica ($\kappa_g$).
		\item Ecuaciones diferenciales de las líneas geodésicas. Coordenadas geodésicas polares y paralelas.
		\item Superficies de curvatura gaussiana constante: esfera ($K > 0$), plano ($K = 0$) y pseudoesfera ($K < 0$).
		\item Teorema global de Gauss-Bonnet y su vínculo con la característica de Euler-Poincaré ($\chi$).
		\item Tópicos adicionales: curvas paralelas, superficies regladas y desarrollables, dinámica de partículas confinadas a superficies, y formulación de tubos envolventes en $\mathbb{R}^2$ y $\mathbb{R}^3$.
	\end{itemize}
	
	\subsubsection*{Aplicaciones Computacionales en el Mundo Real}
	\begin{itemize}
		\item \textbf{Navegación de Drones y Rovers en Terrenos Irregulares:} Algoritmos de distancia mínima y trayectorias geodésicas numéricas sobre modelos digitales de elevación (\emph{Fast Marching Method}).
		\item \textbf{Manufactura con Materiales Rígidos:} Construcción de paneles de fuselaje aeronáutico y patrones textiles restringidos por superficies desarrollables ($K = 0$).
		\item \textbf{Modelado Vascular y Redes de Tuberías:} Generación algorítmica de conductos tubulares 3D a lo largo de trayectorias esqueléticas extraídas de angiotomografías.
		\item \textbf{Auditoría Topológica de Mallas:} Aplicación del Teorema de Gauss-Bonnet discreto para verificar la validez de modelos cerrados para manufactura aditiva.
	\end{itemize}
	
	\subsubsection*{Actividades y Evaluaciones en Moodle}
	\begin{itemize}
		\item \textbf{Laboratorio Computacional 4 (Jupyter Notebook):} Integración numérica del sistema de ecuaciones diferenciales no lineales de las geodésicas sobre superficies de revolución.
		\item \textbf{Foro de Investigación Avanzada:} Revisión de artículos científicos de conferencias ACM SIGGRAPH o IEEE CVPR que aplican geodésicas en aprendizaje profundo geométrico (\emph{Geometric Deep Learning}).
	\end{itemize}
	
	\section{Metodología de Enseñanza en Entorno Virtual (Moodle)}
	
	La dinámica del curso está organizada para el trabajo autónomo guiado y la colaboración académica rigurosa:
	
	\begin{itemize}
		\item \textbf{Lecciones Interactivas y Lecturas Dirigidas:} Disponibilidad asincrónica de notas teóricas redactadas con soporte matemático MathJax y animaciones tridimensionales de apoyo.
		\item \textbf{Laboratorios con Autoevaluación:} Cuadernos interactivos en Python provistos con pruebas unitarias para que el estudiante valide sus algoritmos antes de la entrega formal.
		\item \textbf{Talleres de Coevaluación (Moodle Workshop):} Cada estudiante revisa de manera anónima y fundamentada las soluciones y el código de al menos dos pares.
		\item \textbf{Seminarios Sincrónicos de Asesoría:} Sesiones quincenales por videoconferencia para despejar dudas complejas sobre deducciones analíticas y estabilidad computacional.
	\end{itemize}
	
	\section{Esquema de Evaluación Continua}
	
	\begin{center}
		\begin{tabularx}{\textwidth}{lXr}
			\toprule
			\textbf{Instrumento de Evaluación} & \textbf{Herramienta en Moodle} & \textbf{Ponderación} \\
			\midrule
			4 Laboratorios Computacionales & Módulo de Tareas (Jupyter Notebooks) & 35\% \\
			2 Talleres de Coevaluación por Pares & Módulo Taller (Moodle Workshop) & 15\% \\
			4 Pruebas Analítico-Simbólicas & Cuestionarios STACK / Formulas & 15\% \\
			Participación en Foros Especializados & Foros de Discusión Académica & 10\% \\
			Proyecto Final Integrador Aplicado & Entrega de Monografía y Defensa Oral & 25\% \\
			\midrule
			\textbf{Calificación Total} & & \textbf{100\%} \\
			\bottomrule
		\end{tabularx}
	\end{center}
	
	\section{Proyecto Final Integrador (Líneas de Investigación)}
	
	El estudiante elegirá una línea de investigación aplicada para redactar una monografía de 6 a 8 páginas en formato estándar de artículo de investigación y defender sus resultados en sesión sincrónica:
	
	\begin{enumerate}
		\item \textbf{Reconstrucción de Superficies por Mínimos Cuadrados Móviles (MLS):} Estimación local del plano tangente y curvaturas principales a partir de nubes de puntos con ruido sensorial.
		\item \textbf{Planificación Geodésica de Robots Móviles:} Optimización de trayectorias sobre topografías irregulares minimizando la curvatura geodésica y el gasto energético.
		\item \textbf{Generación de Curvas Suaves para Mecanizado en 5 Ejes:} Interpolación $C^2$ con control de torsión a partir de trayectorias discretas en código G.
		\item \textbf{Análisis Geométrico y Topológico de Estructuras Anatómicas 3D:} Implementación del operador de Laplace-Beltrami y el teorema de Gauss-Bonnet sobre mallas biomédicas reales.
	\end{enumerate}
	

	
\end{document}